% ---
% id: foundations/req:declared-order
% kind: req
% title: "Summation order declared with the data"
% summary: >-
%   Where a predicate on declared data is a sum, its order of accumulation is
%   declared with the alphabet that indexes it and travels with the data. Every
%   consumer, and every declarer, that decides the predicate decides it in that
%   order.
% prominence: standard
% status: accepted
% uses: [def:declared-order, def:declared-data]
% assumes: []
% refs: []
% provenance: newly-derived
% lean:
% ---

\begin{requirement}\label{req:declared-order}
Where a predicate on declared data is a sum, the order in which the sum is
accumulated is part of the declared data: it is declared with the alphabet
that indexes the sum, as that alphabet's declared order
(Definition~\ref{def:declared-order}), or with an alphabet containing it,
whose declared order it then carries by restriction; and it travels with the
artefact that carries the data. Where a declared order arrives both with the
alphabet and with an alphabet containing it, the two are one order: the first
is the restriction of the second. A consumer that receives two such orders
that differ takes the order of the containing alphabet, decides the predicate
in its restriction, and refuses nothing for the difference. Every consumer that decides the predicate,
and every declarer that decides it of data it is about to declare
(Definition~\ref{def:declared-data}), decides it in that order, and no
consumer substitutes an order of its own; a supplier deciding its own
conditions at slacks declared to it is not bound by this. Only where no
declared order arrives with the alphabet or with an alphabet containing it is
the declaration's own listing of the alphabet taken as its declared order.
Where such a predicate reads the trace of a declared matrix, the trace is
accumulated in increasing row index.
\end{requirement}
